%% ============================================================
%%  Chapter 7 — Testing Strategy and Configuration
%%  07_testing_strategy.tex
%% ============================================================
\chapter{Testing Strategy and Configuration}
\label{chap:testing}

\chapterintro{%
  This chapter documents the testing philosophy, infrastructure, and
  configuration of the Clinic Management API test suite. It covers the
  distinction between unit and feature tests in the context of this
  project, explains why a dedicated MySQL test database is used instead
  of an in-memory SQLite substitute, documents the \texttt{phpunit.xml}
  configuration file, describes the \texttt{RefreshDatabase} trait and
  the Passport \texttt{actingAs()} authentication helper, and provides
  the full source for the base \texttt{TestCase} class and the model
  factories that underpin every test.
}

%% ─────────────────────────────────────────────────────────────────────────
\section{Testing Philosophy}
\label{sec:testing-philosophy}
%% ─────────────────────────────────────────────────────────────────────────

The test suite for the Clinic Management API is built around a single
guiding principle: \textbf{tests must exercise the same code path that
production requests follow}. This principle has one significant
architectural consequence that distinguishes this project from many
Laravel applications: the tests run against a real MySQL database, not
an in-memory SQLite substitute.

This choice is deliberate and non-negotiable for three reasons.

\begin{enumerate}
  \item \textbf{FULLTEXT search.} The patient search endpoint uses
    MySQL's \texttt{MATCH ... AGAINST} full-text search syntax via
    Eloquent's \texttt{whereFullText()} method. SQLite does not support
    this syntax. A test that substitutes SQLite would be forced to
    mock the search query entirely, testing nothing of substance.

  \item \textbf{TIME column formatting.} MySQL returns \texttt{TIME}
    column values over the wire as \texttt{HH:MM:SS} strings. The
    application contains explicit normalisation logic
    (\texttt{substr(\$t, 0, 5)}) that converts this to \texttt{HH:MM}
    before slot comparison. SQLite represents time values differently;
    a test database that does not reproduce MySQL's wire format cannot
    catch regressions in this normalisation.

  \item \textbf{Unique composite indexes.} The double-booking prevention
    backstop relies on a MySQL unique composite index on
    \texttt{(doctor\_id, appointment\_date, start\_time)}. SQLite
    handles unique constraint violations through a different code path
    and with different exception types, which would make the test for
    the database-level backstop unreliable.
\end{enumerate}

The test suite therefore mandates a dedicated MySQL database named
\texttt{clinic\_management\_test}, configured separately from the
production database. This database is created once and its schema is
refreshed per-test via the \texttt{RefreshDatabase} trait, which wraps
each test in a database transaction that is rolled back after the test
completes.

%% ─────────────────────────────────────────────────────────────────────────
\section{Unit Tests vs.\ Feature Tests}
\label{sec:unit-vs-feature}
%% ─────────────────────────────────────────────────────────────────────────

PHPUnit categorises tests into test suites, and Laravel conventionally
maps these to two directories: \texttt{tests/Unit/} and
\texttt{tests/Feature/}. The distinction is meaningful in this project.

\subsection{Unit Tests}

A unit test exercises a single class or function in complete isolation.
All dependencies are replaced with mocks or stubs. Unit tests do not
touch the database, the HTTP stack, or the filesystem.

In the Clinic Management API, unit tests are appropriate for:

\begin{itemize}
  \item \textbf{Enum behaviour.} Testing that
    \texttt{AppointmentStatus::PENDING->canTransitionTo(AppointmentStatus::CONFIRMED)}
    returns \texttt{true} and that
    \texttt{AppointmentStatus::COMPLETED->canTransitionTo(AppointmentStatus::PENDING)}
    returns \texttt{false}.
  \item \textbf{Slot generation algorithm.} Testing the pure arithmetic
    logic of the slot generator given a fixed start time, end time, and
    slot duration, without any database involvement.
  \item \textbf{Exception message formatting.} Testing that
    \phpclass{DuplicateScheduleException} constructs the correct message
    string for each day-of-week integer.
\end{itemize}

\subsection{Feature Tests}

A feature test exercises the complete HTTP request-to-response pipeline:
routing, middleware, form request validation, controllers, services,
repositories, Eloquent models, and the database. Every test in the
primary test suite is a feature test.

Feature tests are written using Laravel's \texttt{TestCase} (which
extends PHPUnit's \texttt{TestCase}) and use the HTTP testing helpers
(\texttt{getJson()}, \texttt{postJson()}, etc.) to simulate real HTTP
requests. The response object is then asserted against for status code,
JSON structure, and specific field values.

\begin{table}[h!]
  \centering
  \caption{Comparison of unit and feature test characteristics.}
  \label{tab:unit-vs-feature}
  \renewcommand{\arraystretch}{1.5}
  \begin{tabularx}{\linewidth}{@{}L{3.5cm} X X@{}}
    \toprule
    \rowcolor{PrimaryBlue}
    \tblhdr{Characteristic} & \tblhdr{Unit Test} & \tblhdr{Feature Test} \\
    \midrule
    Scope
      & Single class or method
      & Full HTTP stack \\
    \rowalt
    Database
      & None (mocked or stubbed)
      & Real MySQL test database \\
    Dependencies
      & Mocked via PHPUnit mocks
      & Real instances via Laravel DI container \\
    \rowalt
    Speed
      & Very fast (milliseconds)
      & Slower (tens of milliseconds per test) \\
    Confidence
      & Low (integration gaps uncovered)
      & High (exercises production code path) \\
    \rowalt
    Primary use in this project
      & Enum logic, pure algorithms
      & All business scenarios \\
    \bottomrule
  \end{tabularx}
\end{table}

\begin{notebox}[This Project's Emphasis]
  The overwhelming majority of the \testcount{} tests in this suite are
  feature tests. The decision to invest in feature tests over unit tests
  reflects the architecture of the application: the business logic lives
  in service classes that are deeply integrated with Eloquent and the
  database. Mocking Eloquent to unit-test a service class produces tests
  that verify the mock, not the service.
\end{notebox}

%% ─────────────────────────────────────────────────────────────────────────
\section{PHPUnit Configuration}
\label{sec:phpunit-config}
%% ─────────────────────────────────────────────────────────────────────────

\subsection{The phpunit.xml File}

PHPUnit reads its configuration from \texttt{phpunit.xml} in the
project root. The configuration defines the test suites, the code
coverage settings, and the environment variables that override
\texttt{.env} values during the test run.

\begin{lstlisting}[style=xml, caption={\texttt{phpunit.xml} --- complete
  PHPUnit 11 configuration for the Clinic Management API.}]
<?xml version="1.0" encoding="UTF-8"?>
<phpunit xmlns:xsi="http://www.w3.org/2001/XMLSchema-instance"
         xsi:noNamespaceSchemaLocation="vendor/phpunit/phpunit/phpunit.xsd"
         bootstrap="vendor/autoload.php"
         colors="true"
         stopOnFailure="false"
         processIsolation="false"
         backupGlobals="false"
>
  <!--
    ┌─────────────────────────────────────────────────────────────────┐
    │  Test Suites                                                    │
    └─────────────────────────────────────────────────────────────────┘
  -->
  <testsuites>
    <testsuite name="Unit">
      <directory>tests/Unit</directory>
    </testsuite>
    <testsuite name="Feature">
      <directory>tests/Feature</directory>
    </testsuite>
  </testsuites>

  <!--
    ┌─────────────────────────────────────────────────────────────────┐
    │  Code Coverage Source                                           │
    └─────────────────────────────────────────────────────────────────┘
    Restrict coverage collection to application source only.
    Vendor code and bootstrap files are excluded.
  -->
  <source>
    <include>
      <directory>app</directory>
    </include>
    <exclude>
      <directory>app/Console</directory>
    </exclude>
  </source>

  <!--
    ┌─────────────────────────────────────────────────────────────────┐
    │  Environment Variables                                          │
    └─────────────────────────────────────────────────────────────────┘
    These values override the .env file for the duration of the test
    run. They point all database activity at the dedicated test
    database and disable external side-effects.
  -->
  <php>
    <!-- Application environment -->
    <env name="APP_ENV"   value="testing"/>
    <env name="APP_DEBUG" value="true"/>

    <!--
      Database — dedicated MySQL test schema.
      SQLite :memory: is NOT used because the application depends on
      MySQL-specific features: FULLTEXT indexes (patient search),
      TIME column wire format (slot normalisation), and composite
      unique constraints (double-booking backstop).
    -->
    <env name="DB_CONNECTION" value="mysql"/>
    <env name="DB_HOST"       value="127.0.0.1"/>
    <env name="DB_PORT"       value="3306"/>
    <env name="DB_DATABASE"   value="clinic_management_test"/>
    <env name="DB_USERNAME"   value="root"/>
    <env name="DB_PASSWORD"   value=""/>

    <!-- Cache — array driver so no Redis is required in CI -->
    <env name="CACHE_DRIVER"   value="array"/>

    <!-- Queue — sync driver prevents async jobs from running
         asynchronously during tests, ensuring jobs complete
         before assertions are evaluated -->
    <env name="QUEUE_CONNECTION" value="sync"/>

    <!-- Session — array driver (no cookie handling needed) -->
    <env name="SESSION_DRIVER" value="array"/>

    <!-- Mail — log driver captures outgoing mail without
         sending, preventing test emails reaching real inboxes -->
    <env name="MAIL_MAILER" value="log"/>

    <!-- Passport — use a pre-generated testing key pair so that
         passport:install does not need to run in CI -->
    <env name="PASSPORT_PERSONAL_ACCESS_CLIENT_ID"     value="1"/>
    <env name="PASSPORT_PERSONAL_ACCESS_CLIENT_SECRET" value="testing-secret"/>
  </php>

</phpunit>
\end{lstlisting}

\subsection{Why Not SQLite In-Memory?}

Many Laravel tutorials recommend the \texttt{DB\_CONNECTION=sqlite} and
\texttt{DB\_DATABASE=:memory:} combination for test isolation. This
approach offers genuine speed advantages --- SQLite does not require a
running database server and its in-memory mode creates and destroys the
schema instantaneously. However, it is incompatible with this application
for the reasons enumerated in Section~\ref{sec:testing-philosophy}.

\begin{warningbox}[Do Not Substitute SQLite]
  If \texttt{DB\_CONNECTION} is changed to \texttt{sqlite} and
  \texttt{DB\_DATABASE} is changed to \texttt{:memory:} in
  \texttt{phpunit.xml}, the following categories of tests will fail or
  produce false positives:
  \begin{itemize}
    \item Patient full-text search tests --- SQLite has no
      \texttt{MATCH ... AGAINST} syntax.
    \item Slot availability tests --- SQLite's time representation
      differs from MySQL's \texttt{HH:MM:SS} wire format.
    \item Double-booking database backstop tests --- SQLite raises a
      different PDO exception for unique constraint violations.
  \end{itemize}
\end{warningbox}

\subsection{Creating the Test Database}

The test database must be created manually before the first test run.
It is never created automatically by the test suite.

\begin{lstlisting}[style=bash, caption={One-time setup: create the
  dedicated MySQL test database.}]
# Connect to MySQL and create the schema
mysql -u root -p -e "CREATE DATABASE IF NOT EXISTS clinic_management_test \
  CHARACTER SET utf8mb4 COLLATE utf8mb4_unicode_ci;"

# Run all migrations against the test database
php artisan migrate --env=testing

# (Optional) Confirm the schema is present
php artisan migrate:status --env=testing
\end{lstlisting}

Because the \texttt{RefreshDatabase} trait wraps each individual test
in a transaction that is rolled back at the end, the schema itself
persists across the entire test run. The \texttt{migrate} command above
only needs to be run once (or after a new migration is added).

%% ─────────────────────────────────────────────────────────────────────────
\section{The RefreshDatabase Trait}
\label{sec:refresh-database}
%% ─────────────────────────────────────────────────────────────────────────

The \texttt{Illuminate\textbackslash{}Foundation\textbackslash{}Testing\textbackslash{}RefreshDatabase}
trait is applied to the base \phpclass{TestCase} class, meaning every
feature test in the suite inherits it automatically.

\subsection{How RefreshDatabase Works}

When \texttt{RefreshDatabase} is active:

\begin{enumerate}
  \item \textbf{Before the first test in the class}, the trait checks
    whether the database schema needs to be refreshed. If the schema
    is out of date (i.e.\ unapplied migrations exist), it runs
    \texttt{php artisan migrate --env=testing} to bring the schema
    up to date.

  \item \textbf{Before each individual test method}, the trait begins
    a database transaction (\texttt{START TRANSACTION}).

  \item \textbf{The test executes}, reading and writing database records
    freely. Any factory-generated records, API-created resources, or
    explicit \texttt{DB::table()->insert()} calls land inside this
    transaction.

  \item \textbf{After each individual test method}, the trait rolls back
    the transaction (\texttt{ROLLBACK}), returning the database to the
    exact state it was in before the test started.
\end{enumerate}

The result is complete test isolation: no test can observe the side
effects of any other test. The rollback strategy is significantly faster
than truncating all tables between tests because it avoids full table
scans and re-seeding.

\begin{notebox}[Transaction Nesting]
  The \texttt{RefreshDatabase} rollback works because each test's
  database writes are nested inside the trait's outer transaction. If
  application code explicitly calls \texttt{DB::beginTransaction()} and
  \texttt{DB::commit()}, the inner commit does not escape the outer
  transaction. This means that the \texttt{PrescriptionService}'s use
  of \texttt{DB::transaction()} for atomic prescription creation is
  fully compatible with \texttt{RefreshDatabase}.
\end{notebox}

\subsection{DatabaseMigrations vs.\ RefreshDatabase}

Laravel provides two traits for database isolation:
\texttt{RefreshDatabase} and \texttt{DatabaseMigrations}. The
difference is significant:

\begin{table}[h!]
  \centering
  \caption{Comparison of \texttt{RefreshDatabase} and
    \texttt{DatabaseMigrations} traits.}
  \label{tab:db-traits}
  \renewcommand{\arraystretch}{1.5}
  \begin{tabularx}{\linewidth}{@{}L{4cm} X X@{}}
    \toprule
    \rowcolor{PrimaryBlue}
    \tblhdr{Trait} & \tblhdr{Isolation Mechanism} & \tblhdr{Speed} \\
    \midrule
    \texttt{RefreshDatabase}
      & Wraps each test in a database transaction; rolls back after test
      & Fast --- no table truncation or re-migration \\
    \rowalt
    \texttt{DatabaseMigrations}
      & Rolls back all migrations after each test class, then re-runs
        them before the next
      & Slow --- full migrate down/up cycle per class \\
    \bottomrule
  \end{tabularx}
\end{table}

This project uses \texttt{RefreshDatabase} exclusively.

%% ─────────────────────────────────────────────────────────────────────────
\section{Passport Authentication in Tests}
\label{sec:passport-testing}
%% ─────────────────────────────────────────────────────────────────────────

Authenticating requests in feature tests requires special handling
because Laravel Passport's token verification pipeline normally reads
from the \texttt{oauth\_access\_tokens} table and validates RSA
signatures. Two strategies are available.

\subsection{Strategy 1: actingAs() with the api Guard}

Laravel's \texttt{actingAs(\$user, 'api')} helper bypasses the full
Passport token verification pipeline and directly sets the authenticated
user on the request for the duration of the test. This is the
recommended approach for the vast majority of tests because it is fast,
requires no token issuance, and keeps the test focused on the business
behaviour being tested rather than the authentication mechanics.

\begin{lstlisting}[style=php, caption={Using \texttt{actingAs()} to
  authenticate as a specific user in a feature test.}]
public function test_doctor_can_list_their_own_appointments(): void
{
    $clinic  = Clinic::factory()->create();
    $doctor  = User::factory()->doctor()->forClinic($clinic)->create();

    // Create 3 appointments for this doctor
    Appointment::factory()->count(3)->for($doctor, 'doctor')->create([
        'clinic_id' => $clinic->id,
    ]);

    // Create 2 appointments for a different doctor — should not appear
    $otherDoctor = User::factory()->doctor()->forClinic($clinic)->create();
    Appointment::factory()->count(2)->for($otherDoctor, 'doctor')->create([
        'clinic_id' => $clinic->id,
    ]);

    $response = $this->actingAs($doctor, 'api')
                     ->getJson('/api/v1/doctor/appointments');

    $response->assertOk()
             ->assertJsonCount(3, 'data')
             ->assertJsonPath('success', true);
}
\end{lstlisting}

\subsection{Strategy 2: Real Token Issuance}

For tests that specifically verify the login flow, token revocation, or
the Passport token response format, a real token is issued by calling
the login endpoint directly. This requires the Passport client records
to be seeded in the test database.

The \texttt{Passport::actingAs()} method (distinct from Laravel's
\texttt{actingAs()}) can also be used to create a token in memory and
attach it to the request without a database write, which is useful for
testing token-scope behaviour.

\begin{lstlisting}[style=php, caption={Testing the login endpoint with a
  real credential round-trip.}]
public function test_valid_credentials_return_access_token(): void
{
    $user = User::factory()->create([
        'email'    => 'test@example.com',
        'password' => bcrypt('password123'),
        'role'     => UserRole::DOCTOR,
    ]);

    $response = $this->postJson('/api/v1/auth/login', [
        'email'    => 'test@example.com',
        'password' => 'password123',
    ]);

    $response->assertOk()
             ->assertJsonStructure([
                 'success',
                 'message',
                 'data' => [
                     'user',
                     'access_token',
                     'token_type',
                 ],
             ])
             ->assertJsonPath('data.token_type', 'Bearer')
             ->assertJsonPath('success', true);
}
\end{lstlisting}

\subsection{Passport::actingAs() for Scope Testing}

When the test needs to verify that a specific OAuth2 scope is required,
\texttt{Passport::actingAs()} allows a token with a specific scope
array to be injected without going through the issuance pipeline.
This project does not use OAuth2 scopes (Personal Access Tokens carry
an empty scope array), so this pattern is included for completeness
rather than active use.

\begin{lstlisting}[style=php, caption={\texttt{Passport::actingAs()}
  --- scope-aware authentication for hypothetical scope-gated endpoints.}]
use Laravel\Passport\Passport;

public function test_token_with_insufficient_scope_is_rejected(): void
{
    $user = User::factory()->doctor()->create();

    // Issue an in-memory token with no scopes
    Passport::actingAs($user, []);

    $response = $this->getJson('/api/v1/doctor/appointments');

    // In the current implementation all PATs have full access,
    // so this would return 200. Shown for illustration only.
    $response->assertOk();
}
\end{lstlisting}

%% ─────────────────────────────────────────────────────────────────────────
\section{The Base TestCase Class}
\label{sec:base-testcase}
%% ─────────────────────────────────────────────────────────────────────────

The base \texttt{TestCase} class at \texttt{tests/TestCase.php} is
the foundation from which every test class inherits. It applies shared
traits, registers Passport keys, seeds the Passport client records
required for token issuance tests, and provides shared helper methods
used across the test suite.

\begin{lstlisting}[style=php, caption={\texttt{tests/TestCase.php} ---
  Base test case with Passport configuration and shared helpers.}]
<?php

namespace Tests;

use App\Models\Clinic;
use App\Models\User;
use Illuminate\Foundation\Testing\RefreshDatabase;
use Illuminate\Foundation\Testing\TestCase as BaseTestCase;
use Laravel\Passport\Passport;

abstract class TestCase extends BaseTestCase
{
    use RefreshDatabase;

    /**
     * Run before each test method.
     *
     * Passport::loadKeysFrom() tells Passport where to find the RSA
     * keys without requiring passport:install to have been executed.
     * In the testing environment the keys are loaded from the
     * storage/ directory where they were generated once during initial
     * project setup.
     */
    protected function setUp(): void
    {
        parent::setUp();

        // Point Passport at the pre-generated key files.
        // The testing environment does not run passport:install; the
        // keys generated during local setup are committed to storage/
        // (but never to version control in a real project).
        Passport::loadKeysFrom(storage_path());
    }

    // ──────────────────────────────────────────────────────────────────
    // Convenience Authentication Helpers
    // ──────────────────────────────────────────────────────────────────

    /**
     * Authenticate the test client as a Super Administrator.
     */
    protected function actingAsSuperAdmin(?User $user = null): static
    {
        $user ??= User::factory()->superAdmin()->create();
        return $this->actingAs($user, 'api');
    }

    /**
     * Authenticate the test client as a Doctor bound to the given clinic.
     */
    protected function actingAsDoctor(Clinic $clinic, ?User $user = null): static
    {
        $user ??= User::factory()->doctor()->forClinic($clinic)->create();
        return $this->actingAs($user, 'api');
    }

    /**
     * Authenticate the test client as an Assistant bound to the given clinic.
     */
    protected function actingAsAssistant(Clinic $clinic, ?User $user = null): static
    {
        $user ??= User::factory()->assistant()->forClinic($clinic)->create();
        return $this->actingAs($user, 'api');
    }

    // ──────────────────────────────────────────────────────────────────
    // JSON Structure Helpers
    // ──────────────────────────────────────────────────────────────────

    /**
     * Assert that the response carries the standard success envelope
     * and the expected HTTP status code.
     */
    protected function assertSuccessResponse(
        \Illuminate\Testing\TestResponse $response,
        int $status = 200,
    ): void {
        $response->assertStatus($status)
                 ->assertJsonPath('success', true)
                 ->assertJsonStructure(['success', 'message', 'data']);
    }

    /**
     * Assert that the response carries the standard error envelope
     * and the expected HTTP status code.
     */
    protected function assertErrorResponse(
        \Illuminate\Testing\TestResponse $response,
        int    $status,
        string $message,
    ): void {
        $response->assertStatus($status)
                 ->assertJsonPath('success', false)
                 ->assertJsonPath('message', $message);
    }
}
\end{lstlisting}

\begin{notebox}[RefreshDatabase on the Base Class]
  Applying \texttt{RefreshDatabase} to the base \texttt{TestCase} rather
  than to each individual test class ensures that no test class can
  accidentally omit it and pollute the database state for subsequent
  tests. Every test that extends this \phpclass{TestCase} is
  automatically wrapped in a transaction.
\end{notebox}

%% ─────────────────────────────────────────────────────────────────────────
\section{Model Factories}
\label{sec:factories}
%% ─────────────────────────────────────────────────────────────────────────

All test data is generated through Laravel model factories. The
factories reside in \texttt{database/factories/} and use the Faker
library to produce realistic, randomised values. Each factory defines
state methods that compose to build complex test scenarios with minimal
boilerplate.

\subsection{UserFactory}

The \phpclass{UserFactory} is the most frequently used factory in the
suite. Its state methods (\texttt{superAdmin()}, \texttt{doctor()},
\texttt{assistant()}, \texttt{forClinic()}, \texttt{inactive()}) allow
any combination of role and clinic to be expressed in a single fluent
chain.

\begin{lstlisting}[style=php, caption={\texttt{database/factories/UserFactory.php}
  --- Role and clinic state methods for composable test data.}]
<?php

namespace Database\Factories;

use App\Enums\UserRole;
use App\Models\Clinic;
use Illuminate\Database\Eloquent\Factories\Factory;
use Illuminate\Support\Facades\Hash;

class UserFactory extends Factory
{
    public function definition(): array
    {
        return [
            'name'      => fake()->name(),
            'email'     => fake()->unique()->safeEmail(),
            'password'  => Hash::make('password'),
            'role'      => UserRole::DOCTOR,
            'clinic_id' => Clinic::factory(),
            'is_active' => true,
        ];
    }

    public function superAdmin(): static
    {
        return $this->state(fn(array $attrs) => [
            'role'      => UserRole::SUPER_ADMIN,
            'clinic_id' => null,
        ]);
    }

    public function doctor(): static
    {
        return $this->state(fn(array $attrs) => [
            'role' => UserRole::DOCTOR,
        ]);
    }

    public function assistant(): static
    {
        return $this->state(fn(array $attrs) => [
            'role' => UserRole::ASSISTANT,
        ]);
    }

    public function forClinic(Clinic $clinic): static
    {
        return $this->state(fn(array $attrs) => [
            'clinic_id' => $clinic->id,
        ]);
    }

    public function inactive(): static
    {
        return $this->state(fn(array $attrs) => [
            'is_active' => false,
        ]);
    }
}
\end{lstlisting}

\subsection{ClinicFactory}

\begin{lstlisting}[style=php, caption={\texttt{database/factories/ClinicFactory.php}.}]
<?php

namespace Database\Factories;

use Illuminate\Database\Eloquent\Factories\Factory;

class ClinicFactory extends Factory
{
    public function definition(): array
    {
        return [
            'name'      => fake()->unique()->company() . ' Clinic',
            'email'     => fake()->unique()->companyEmail(),
            'phone'     => fake()->phoneNumber(),
            'address'   => fake()->address(),
            'is_active' => true,
        ];
    }

    public function inactive(): static
    {
        return $this->state(fn(array $attrs) => [
            'is_active' => false,
        ]);
    }
}
\end{lstlisting}

\subsection{PatientFactory}

\begin{lstlisting}[style=php, caption={\texttt{database/factories/PatientFactory.php}.}]
<?php

namespace Database\Factories;

use App\Enums\Gender;
use App\Models\Clinic;
use Illuminate\Database\Eloquent\Factories\Factory;

class PatientFactory extends Factory
{
    public function definition(): array
    {
        return [
            'name'          => fake()->name(),
            'phone'         => fake()->unique()->phoneNumber(),
            'email'         => fake()->optional()->safeEmail(),
            'date_of_birth' => fake()->date('Y-m-d', '-18 years'),
            'gender'        => fake()->randomElement(Gender::cases()),
            'address'       => fake()->optional()->address(),
            'clinic_id'     => Clinic::factory(),
        ];
    }

    public function forClinic(Clinic $clinic): static
    {
        return $this->state(fn(array $attrs) => [
            'clinic_id' => $clinic->id,
        ]);
    }

    public function male(): static
    {
        return $this->state(fn(array $attrs) => [
            'gender' => Gender::MALE,
        ]);
    }

    public function female(): static
    {
        return $this->state(fn(array $attrs) => [
            'gender' => Gender::FEMALE,
        ]);
    }
}
\end{lstlisting}

\subsection{DoctorScheduleFactory}

\begin{lstlisting}[style=php, caption={\texttt{database/factories/DoctorScheduleFactory.php}
  --- Day-of-week state methods for readable schedule setup.}]
<?php

namespace Database\Factories;

use App\Models\User;
use Illuminate\Database\Eloquent\Factories\Factory;

class DoctorScheduleFactory extends Factory
{
    public function definition(): array
    {
        return [
            'doctor_id'     => User::factory()->doctor(),
            'clinic_id'     => null, // resolved from doctor's clinic_id in service
            'day_of_week'   => fake()->numberBetween(0, 6),
            'start_time'    => '09:00',
            'end_time'      => '17:00',
            'slot_duration' => 30,
        ];
    }

    public function sunday(): static
    {
        return $this->state(fn(array $attrs) => ['day_of_week' => 0]);
    }

    public function monday(): static
    {
        return $this->state(fn(array $attrs) => ['day_of_week' => 1]);
    }

    public function tuesday(): static
    {
        return $this->state(fn(array $attrs) => ['day_of_week' => 2]);
    }

    public function wednesday(): static
    {
        return $this->state(fn(array $attrs) => ['day_of_week' => 3]);
    }

    public function thursday(): static
    {
        return $this->state(fn(array $attrs) => ['day_of_week' => 4]);
    }

    public function friday(): static
    {
        return $this->state(fn(array $attrs) => ['day_of_week' => 5]);
    }

    public function saturday(): static
    {
        return $this->state(fn(array $attrs) => ['day_of_week' => 6]);
    }

    public function slotDuration(int $minutes): static
    {
        return $this->state(fn(array $attrs) => ['slot_duration' => $minutes]);
    }
}
\end{lstlisting}

\subsection{AppointmentFactory}

\begin{lstlisting}[style=php, caption={\texttt{database/factories/AppointmentFactory.php}
  --- Status state methods for lifecycle scenario testing.}]
<?php

namespace Database\Factories;

use App\Enums\AppointmentStatus;
use App\Models\Clinic;
use App\Models\Patient;
use App\Models\User;
use Illuminate\Database\Eloquent\Factories\Factory;

class AppointmentFactory extends Factory
{
    public function definition(): array
    {
        $clinic = Clinic::factory()->create();

        return [
            'doctor_id'        => User::factory()->doctor()->forClinic($clinic),
            'patient_id'       => Patient::factory()->forClinic($clinic),
            'clinic_id'        => $clinic->id,
            'booked_by'        => User::factory()->assistant()->forClinic($clinic),
            'appointment_date' => fake()->dateTimeBetween('now', '+3 months')
                                        ->format('Y-m-d'),
            'start_time'       => '09:00',
            'end_time'         => '09:30',
            'status'           => AppointmentStatus::PENDING,
            'notes'            => fake()->optional()->sentence(),
        ];
    }

    public function pending(): static
    {
        return $this->state(fn(array $attrs) => [
            'status' => AppointmentStatus::PENDING,
        ]);
    }

    public function confirmed(): static
    {
        return $this->state(fn(array $attrs) => [
            'status' => AppointmentStatus::CONFIRMED,
        ]);
    }

    public function completed(): static
    {
        return $this->state(fn(array $attrs) => [
            'status' => AppointmentStatus::COMPLETED,
        ]);
    }

    public function cancelled(): static
    {
        return $this->state(fn(array $attrs) => [
            'status' => AppointmentStatus::CANCELLED,
        ]);
    }

    public function forDate(string $date): static
    {
        return $this->state(fn(array $attrs) => [
            'appointment_date' => $date,
        ]);
    }

    public function atTime(string $startTime, string $endTime): static
    {
        return $this->state(fn(array $attrs) => [
            'start_time' => $startTime,
            'end_time'   => $endTime,
        ]);
    }
}
\end{lstlisting}

\subsection{PrescriptionFactory}

\begin{lstlisting}[style=php, caption={\texttt{database/factories/PrescriptionFactory.php}.}]
<?php

namespace Database\Factories;

use App\Models\Appointment;
use Illuminate\Database\Eloquent\Factories\Factory;

class PrescriptionFactory extends Factory
{
    public function definition(): array
    {
        return [
            'appointment_id' => Appointment::factory()->completed(),
            'diagnosis'      => fake()->sentence(8),
            'notes'          => fake()->optional()->paragraph(),
        ];
    }
}
\end{lstlisting}

\subsection{PrescriptionItemFactory}

\begin{lstlisting}[style=php, caption={\texttt{database/factories/PrescriptionItemFactory.php}.}]
<?php

namespace Database\Factories;

use App\Models\Prescription;
use Illuminate\Database\Eloquent\Factories\Factory;

class PrescriptionItemFactory extends Factory
{
    public function definition(): array
    {
        $medicines = [
            'Ibuprofen', 'Paracetamol', 'Amoxicillin', 'Omeprazole',
            'Metformin', 'Atorvastatin', 'Amlodipine', 'Diclofenac',
        ];

        return [
            'prescription_id' => Prescription::factory(),
            'medicine_name'   => fake()->randomElement($medicines),
            'dosage'          => fake()->randomElement([
                '500mg', '250mg', '1g', '400mg', '200mg', '100mg',
            ]),
            'frequency'       => fake()->randomElement([
                'Once daily', 'Twice daily', 'Three times daily',
                'Every 8 hours', 'As needed',
            ]),
            'duration'        => fake()->randomElement([
                '3 days', '5 days', '7 days', '10 days', '14 days',
            ]),
            'notes'           => fake()->optional()->sentence(),
        ];
    }
}
\end{lstlisting}

%% ─────────────────────────────────────────────────────────────────────────
\section{Test Directory Structure}
\label{sec:test-structure}
%% ─────────────────────────────────────────────────────────────────────────

The test suite is organised by role surface, mirroring the route group
structure of the application. Each role's test classes cover all
endpoints under that role's prefix.

\begin{lstlisting}[style=bash, caption={Test directory tree.}]
tests/
├── TestCase.php                    # Base class with shared helpers
│
├── Unit/
│   ├── Enums/
│   │   ├── AppointmentStatusTest.php  # canTransitionTo() matrix
│   │   └── UserRoleTest.php           # isClinicalStaff(), label()
│   └── Services/
│       └── SlotGeneratorTest.php      # Pure slot arithmetic logic
│
└── Feature/
    ├── Auth/
    │   ├── LoginTest.php              # Valid login, bad creds, inactive
    │   ├── LogoutTest.php             # Token revocation
    │   └── MeTest.php                 # Profile retrieval
    │
    ├── Admin/
    │   ├── DashboardTest.php          # KPI aggregate, date filter
    │   ├── Clinics/
    │   │   ├── ListClinicsTest.php
    │   │   ├── CreateClinicTest.php
    │   │   ├── ShowClinicTest.php
    │   │   ├── UpdateClinicTest.php
    │   │   └── DeleteClinicTest.php
    │   └── Users/
    │       ├── ListUsersTest.php
    │       ├── CreateUserTest.php
    │       ├── ShowUserTest.php
    │       ├── UpdateUserTest.php
    │       ├── DeleteUserTest.php
    │       ├── UpdateUserRoleTest.php  # Self-demotion guard
    │       └── UpdateUserStatusTest.php
    │
    ├── Doctor/
    │   ├── Schedules/
    │   │   ├── ListSchedulesTest.php
    │   │   ├── CreateScheduleTest.php  # Duplicate day guard
    │   │   ├── ShowScheduleTest.php
    │   │   ├── UpdateScheduleTest.php
    │   │   └── DeleteScheduleTest.php
    │   ├── Appointments/
    │   │   ├── ListAppointmentsTest.php # Doctor-scoped listing
    │   │   ├── ShowAppointmentTest.php
    │   │   └── UpdateStatusTest.php     # All transition paths
    │   ├── Patients/
    │   │   ├── ListPatientsTest.php     # FULLTEXT search
    │   │   └── ShowPatientTest.php      # Cross-clinic guard
    │   └── Prescriptions/
    │       ├── CreatePrescriptionTest.php  # Ownership + state guards
    │       ├── ShowPrescriptionTest.php
    │       └── UpdatePrescriptionTest.php
    │
    └── Assistant/
        ├── AvailableSlotsTest.php          # Slot computation
        ├── Patients/
        │   ├── ListPatientsTest.php
        │   ├── CreatePatientTest.php       # Phone uniqueness per clinic
        │   ├── ShowPatientTest.php
        │   └── UpdatePatientTest.php
        └── Appointments/
            ├── ListAppointmentsTest.php    # Clinic-wide listing
            ├── ShowAppointmentTest.php
            ├── BookAppointmentTest.php     # Conflict, out-of-schedule
            ├── RescheduleAppointmentTest.php # State guard
            └── CancelAppointmentTest.php
\end{lstlisting}

%% ─────────────────────────────────────────────────────────────────────────
\section{Running the Test Suite}
\label{sec:running-tests}
%% ─────────────────────────────────────────────────────────────────────────

\subsection{Full Suite}

\begin{lstlisting}[style=bash, caption={Run all tests.}]
php artisan test
# or equivalently:
./vendor/bin/phpunit
\end{lstlisting}

\subsection{Specific Suite or Filter}

\begin{lstlisting}[style=bash, caption={Run only the Feature suite, or filter
  to a specific test class or method.}]
# Feature suite only
php artisan test --testsuite=Feature

# Specific test class
php artisan test --filter=BookAppointmentTest

# Specific test method
php artisan test --filter=test_booking_an_already_taken_slot_returns_conflict

# Specific directory
php artisan test tests/Feature/Doctor/
\end{lstlisting}

\subsection{Coverage Report}

\begin{lstlisting}[style=bash, caption={Generate an HTML coverage report
  (requires Xdebug or PCOV).}]
./vendor/bin/phpunit --coverage-html coverage/

# View in browser
open coverage/index.html
\end{lstlisting}

\subsection{Parallel Execution}

Laravel's \texttt{--parallel} flag runs test classes concurrently across
multiple CPU cores using the \texttt{paratest} package. Because each
test is wrapped in its own transaction, parallel execution is safe.

\begin{lstlisting}[style=bash, caption={Run tests in parallel across 4 workers.}]
php artisan test --parallel --processes=4
\end{lstlisting}

\begin{warningbox}[Parallel Execution and the Test Database]
  When running in parallel, each worker process requires its own
  database connection. Laravel's parallel test runner creates separate
  test database copies (e.g.\ \texttt{clinic\_management\_test\_1},
  \texttt{clinic\_management\_test\_2}) automatically when
  \texttt{--parallel} is specified, provided the database user has
  \texttt{CREATE DATABASE} privileges.
\end{warningbox}

%% ─────────────────────────────────────────────────────────────────────────
\section{Test Configuration Summary}
\label{sec:test-config-summary}
%% ─────────────────────────────────────────────────────────────────────────

\begin{table}[h!]
  \centering
  \caption{Summary of the test environment configuration and its
    rationale.}
  \label{tab:test-config-summary}
  \renewcommand{\arraystretch}{1.5}
  \begin{tabularx}{\linewidth}{@{}L{3.8cm} L{3.5cm} X@{}}
    \toprule
    \rowcolor{PrimaryBlue}
    \tblhdr{Configuration Key} & \tblhdr{Test Value} & \tblhdr{Rationale} \\
    \midrule
    \texttt{DB\_CONNECTION}
      & \texttt{mysql}
      & FULLTEXT search and TIME wire format require MySQL. \\
    \rowalt
    \texttt{DB\_DATABASE}
      & \texttt{clinic\_management\_test}
      & Dedicated schema isolates test activity from the development
        database. \\
    \texttt{CACHE\_DRIVER}
      & \texttt{array}
      & In-memory array cache requires no Redis instance in CI;
        \texttt{Cache::lock()} still functions correctly. \\
    \rowalt
    \texttt{QUEUE\_CONNECTION}
      & \texttt{sync}
      & Dispatched jobs execute immediately within the test request
        cycle, so post-dispatch assertions see the job's side effects. \\
    \texttt{SESSION\_DRIVER}
      & \texttt{array}
      & API tests do not use cookie sessions; array driver avoids
        filesystem writes. \\
    \rowalt
    \texttt{MAIL\_MAILER}
      & \texttt{log}
      & Captures outgoing mail in the Laravel log without hitting an
        SMTP server or a real inbox. \\
    \texttt{APP\_ENV}
      & \texttt{testing}
      & Activates Laravel's testing-specific behaviour, including
        the \texttt{RefreshDatabase} transaction strategy. \\
    \rowalt
    Database isolation
      & \texttt{RefreshDatabase} (transaction rollback)
      & Fastest isolation mechanism; compatible with nested
        \texttt{DB::transaction()} calls in service classes. \\
    Authentication
      & \texttt{actingAs(\$user, 'api')}
      & Bypasses Passport token pipeline; keeps test logic focused
        on the business behaviour under test. \\
    \bottomrule
  \end{tabularx}
\end{table}
